\documentclass[11pt,reqno]{amsart}
\usepackage[T1]{fontenc}
\usepackage[utf8]{inputenc}
\usepackage{lmodern}
\usepackage{amsmath,amssymb,amsfonts,amsthm,mathtools,mathrsfs}
\usepackage{enumitem,microtype}
\usepackage[colorlinks=true,linkcolor=blue,citecolor=blue,urlcolor=blue]{hyperref}
\usepackage[nameinlink,capitalize]{cleveref}
\allowdisplaybreaks
\setlength{\emergencystretch}{2em}

\newtheorem{theorem}{Theorem}[section]
\newtheorem{proposition}[theorem]{Proposition}
\newtheorem{lemma}[theorem]{Lemma}
\newtheorem{corollary}[theorem]{Corollary}
\theoremstyle{definition}
\newtheorem{assumption}[theorem]{Assumption}
\theoremstyle{remark}
\newtheorem{remark}[theorem]{Remark}

\newcommand{\A}{\mathcal A}
\newcommand{\B}{\mathcal B}
\newcommand{\C}{\mathbb C}
\newcommand{\R}{\mathbb R}
\newcommand{\T}{\mathbb T}
\newcommand{\Piop}{\Pi}
\newcommand{\Id}{\mathrm{Id}}
\newcommand{\tr}{\operatorname{tr}}
\newcommand{\dist}{\operatorname{dist}}
\newcommand{\norm}[1]{\left\lVert #1\right\rVert}
\newcommand{\abs}[1]{\left\lvert #1\right\rvert}

\title[Pressure, diffusion, and all-frequency response]{Pressure, Physical Diffusion, and All-Frequency Suspension Response for Moving-Cut Dynamics}
\author{Qian Qi}
\address{Peking University, Beijing 100871, China}
\email{qiqian@pku.edu.cn}
\subjclass[2020]{Primary 37A50, 37C30; Secondary 47A55, 60F17}
\keywords{pressure, suspension flow, Green--Kubo formula, diffusion, high-frequency resolvent}
\date{August 28, 2026}

\begin{document}
\begin{abstract}
We convert the all-order moving-cut response of the companion paper into
physical-time pressure, diffusion, and suspension resolvents.  A finite Banach
bundle is stabilized inside one direct-sum space.  Joint displacement and roof
twists have a differentiable simple eigenvalue; its zero-pressure root gives
physical drift and covariance.  A renewal formula with distinct entry and exit
operators gives the low-frequency flow resolvent.

The diffusion tensor is the symmetrized Green--Kubo series and is proved equal
to the pressure Hessian, asymptotic variance, and Gordin martingale bracket.
For the same nonconjugate moving-cut family, a Diophantine branch roof yields a
polynomial all-frequency resolvent bound and parameter derivatives of every
prescribed finite order.  Finally, for the triangular periodic Lorentz gas
with nearest-neighbour distance one we prove the correct open finite-horizon
window $\sqrt3/4<r<1/2$ and its stability under quantitative $C^2$
deformations.
\end{abstract}
\maketitle
\tableofcontents

\section{Introduction}
The physical-time coefficients require a joint displacement--roof pressure,
not merely a collision-map derivative.  High-frequency family response also
requires a quantitative separation mechanism, not fixed-table estimates plus
compactness.  We provide both on the explicit moving-cut system of Paper I.
The spectral ingredients are classical \cite{Kato1995,KellerLiverani1999,
ParryPollicott1990}; the frequency layer is related to
\cite{Dolgopyat1998,ButterleyLiverani2007}.

\section{Stabilizing a Banach bundle}\label{sec:bundle}
Let $\B_a$ have a finite $C^N$ atlas
$G_{\alpha,a}:\B_a\to\B_\alpha$.  Choose smooth functions $\psi_\alpha$
subordinate to the atlas with
\[
 \sum_\alpha\psi_\alpha(a)^2=1.
\]
On $\B_*=\bigoplus_\alpha\B_\alpha$, define
\[
 J_ah=(\psi_\alpha G_{\alpha,a}h)_\alpha,
 \qquad
 R_a(v_\alpha)=\sum_\alpha\psi_\alpha G_{\alpha,a}^{-1}v_\alpha.
\]

\begin{theorem}[Finite-atlas stabilization]\label{thm:bundle}
One has $R_aJ_a=\Id$.  Thus $P_a=J_aR_a$ is a uniformly bounded projection
onto $E_a=J_a\B_a$, and
$\widehat L_a=J_aL_aR_a\in\mathcal L(\B_*)$ has the same parameter
regularity as the fibre family.  Riesz data around an eigenvalue separated
from zero agree with the fibre data on $E_a$.
\end{theorem}

\begin{proof}
The quadratic partition gives $R_aJ_a=\sum\psi_\alpha^2\Id=\Id$.
All other statements are finite product rules.  The added complementary
spectrum is at zero and does not meet a contour around one.
\end{proof}

\section{Joint pressure and physical time}\label{sec:pressure}
Let $\kappa_a:M\to\R^d$ be a displacement and
$\tau_a:M\to(0,\infty)$ a roof.  Define
\[
 L_{a,q,s}h=L_a(e^{q\cdot\kappa_a-s\tau_a}h).
\]
Assume the graded $N$th-order response theorem of Paper I, multiplier bounds,
and a uniform contour around a simple eigenvalue $\lambda(a,q,s)$ near one.
Put $\mathscr P=\log\lambda$.

\begin{theorem}[Pressure jet]\label{thm:pressure}
The eigenvalue, Riesz projection, reduced resolvent, and pressure are $C^N$ in
$(a,q,s)$.  Every derivative is a finite sum of direct operator letters and
well-typed reduced-resolvent words of total order at most $N$.
\end{theorem}

\begin{proof}
Differentiate the Riesz integral repeatedly using
$D(z-L)^{-1}=(z-L)^{-1}(DL)(z-L)^{-1}$.  The Paper-I ladder makes every word
bounded, and uniform contour separation permits differentiation under the
integral.
\end{proof}

Let $\Lambda_a(q)$ solve
\[
 \mathscr P(a,q,\Lambda_a(q))=0.
\]
Since $\mathscr P_s(a,0,0)=-\bar\tau_a<0$, the implicit-function theorem
applies.

\begin{theorem}[Physical covariance]\label{thm:physical}
Under centered displacement,
\[
 D_q\Lambda_a(0)=0,
 \qquad
 \Sigma_{ij}(a)=\partial_{q_iq_j}\Lambda_a(0)
 ={\mathscr P_{q_iq_j}(a,0,0)\over\bar\tau_a}.
\]
Moreover,
\[
 \partial_a\Sigma_{ij}
 ={\mathscr P_{aq_iq_j}\over\bar\tau}
 -{\mathscr P_{q_iq_j}\partial_a\bar\tau\over\bar\tau^2},
 \qquad D^{\rm phys}=\Sigma/2.
\]
\end{theorem}

\begin{proof}
Differentiate the zero-pressure equation twice in $q$ and once in $a$.
\end{proof}

\section{Oriented renewal}\label{sec:renewal}
For a flow observable $G$, define the entry function
\[
 \mathcal B_{G,a,z}(x)=
 \int_0^{\tau_a(x)}e^{-z(\tau_a(x)-u)}G(x,u)\,du,
\]
and for $F$ the exit functional
\[
 \ell_{F,a,z}(h)=
 \int_M\int_0^{\tau_a(x)}e^{-zu}F(x,u)h(x)\,du\,d\mu_a(x).
\]
Let $L_{a,z}=L_a(e^{-z\tau_a}\cdot)$.

\begin{theorem}[Suspension renewal identity]\label{thm:renewal}
For centered correlations,
\[
 \widehat C_{F,G}(a,z)=H_{a,z}^{\rm local}(F,G)
 +\ell_{F,a,z}(I-L_{a,z})^{-1}\mathcal B_{G,a,z}.
\]
Near zero the pole from the leading eigenvalue cancels, and the centered
quantity has the same parameter regularity as the pressure jet.
\end{theorem}

\begin{proof}
Decompose trajectories by completed returns.  The first and last incomplete
segments give the entry and exit operators; the complete returns sum to the
resolvent.  In
\[
 (I-L_{a,z})^{-1}={\Piop_{a,z}\over1-\lambda(a,0,z)}+R_{a,z}^{\perp},
\]
centering removes the pole and
$\partial_z\lambda(a,0,0)=-\bar\tau_a$ is uniformly nonzero.
\end{proof}

\section{Symmetric covariance}\label{sec:GK}
For centered $\kappa_a$, define
\[
 C_{ij}^{\rm coll}=
 \int\kappa_i\kappa_j\,d\mu+
 \sum_{n\ge1}\int
 (\kappa_i\kappa_j\circ T^n+
  \kappa_j\kappa_i\circ T^n)\,d\mu.
\]

\begin{theorem}[Pressure--Green--Kubo--bracket identity]\label{thm:GK}
Assuming absolute summability,
\[
 \mathscr P_{q_iq_j}=C_{ij}^{\rm coll}
 =\int m_i m_j\,d\mu,
\]
where $m$ is the Gordin martingale increment.  Furthermore,
\[
 \xi^TC^{\rm coll}\xi=
 \lim_{n\to\infty}{1\over n}
 \int\left(\sum_{k=0}^{n-1}\xi\cdot\kappa\circ T^k\right)^2d\mu\ge0.
\]
Hence $C^{\rm coll}$ and $\Sigma=C^{\rm coll}/\bar\tau$ are symmetric
positive semidefinite.  Their kernel is the coboundary subspace in the
declared spectral class.
\end{theorem}

\begin{proof}
Twice differentiate the leading eigenvalue equation.  The two resolvent
insertions expand into the two ordered correlation series.  Expanding the
partial-sum square gives the variance limit; the Poisson decomposition gives
the martingale bracket.
\end{proof}

\begin{corollary}[Square-root modulus]
If $\Sigma(a)\ge\lambda I$, then
\[
 \norm{\Sigma(a)^{1/2}-\Sigma(a')^{1/2}}
 \le{1\over2\sqrt\lambda}\norm{\Sigma(a)-\Sigma(a')}.
\]
\end{corollary}

\begin{proof}
Use the Sylvester equation for the difference of the principal square roots.
\end{proof}

\section{Explicit moving-cut coefficients}\label{sec:explicit}
For the four-branch family of Paper I, choose branch displacements
$c_i\in\R^d$ and fixed positive roofs $\tau_i$.  Then
\[
 \mathscr P(a,q,s)=
 \log\sum_{i=1}^4w_i(a)e^{q\cdot c_i-s\tau_i}.
\]
Let $\bar c(a)=\sum_iw_i(a)c_i$ and
$\widetilde c_i=c_i-\bar c(a)$.

\begin{theorem}[Actual coefficient package]\label{thm:actual}
The pressure and physical root are analytic and
\[
 C^{\rm coll}(a)=
 \sum_iw_i(a)\widetilde c_i(a)\otimes\widetilde c_i(a),
 \qquad
 \Sigma(a)={C^{\rm coll}(a)\over\bar\tau(a)}.
\]
If the centered branch vectors span $\R^d$ for all $a$, then
$\Sigma(a)\ge\lambda I$ uniformly.  The complete low-frequency suspension
response is analytic.
\end{theorem}

\begin{proof}
Paper I's exact innovation theorem makes the fixed-parameter branch symbols
i.i.d. with probabilities $w_i(a)$; all positive-time centered correlations
vanish.  The remaining claims follow from the explicit pressure and compactness
of the smallest covariance eigenvalue.
\end{proof}

\section{Diophantine all-frequency response}\label{sec:HF}
Let
$\omega=(\tau_2-\tau_1,\tau_3-\tau_1,\tau_4-\tau_1)$.

\begin{assumption}[Diophantine roof]\label{ass:dio}
There are $\gamma>0$ and $\nu\ge0$ such that, for $\abs b\ge1$,
\[
 \max_{2\le j\le4}
 \dist(b(\tau_j-\tau_1),2\pi\mathbb Z)
 \ge\gamma(1+\abs b)^{-\nu}.
\]
\end{assumption}

Define
\[
 (L_{a,b}h)(y)=\sum_iw_i(a)e^{-ib\tau_i}
 h(s_{i-1}(a)+w_i(a)y),
 \qquad
 \lambda_a(b)=\sum_iw_i(a)e^{-ib\tau_i}.
\]

\begin{lemma}[Phase gap]\label{lem:phase}
Uniformly in $a$,
\[
 1-\abs{\lambda_a(b)}
 \ge c(1+\abs b)^{-2\nu},
 \qquad\abs b\ge1.
\]
\end{lemma}

\begin{proof}
Use
\[
 1-\abs{\lambda_a(b)}^2
 =2\sum_{i<j}w_iw_j
 [1-\cos(b(\tau_i-\tau_j))]
\]
and the positive minimum branch weight.  The Diophantine separation supplies
one phase difference of size at least
$\gamma(1+\abs b)^{-\nu}$ modulo $2\pi$.
\end{proof}

\begin{theorem}[All-frequency moving-family resolvent]\label{thm:HF}
For every finite $r,k$,
\begin{align}
 \norm{(I-L_{a,b})^{-1}}_{W^{r,1}\to W^{r,1}}
 &\le C_r(1+\abs b)^{2\nu},\label{eq:Rb}\\
 \norm{\partial_a^k(I-L_{a,b})^{-1}}_{W^{r+k,1}\to W^{r,1}}
 &\le C_{r,k}(1+\abs b)^{2\nu(k+1)}.
 \label{eq:Rb-der}
\end{align}
The same estimates hold for $z=\sigma+ib$ with
$0\le\sigma\le\sigma_0$ after enlarging the constants.
\end{theorem}

\begin{proof}
Constants form an invariant block with eigenvalue $\lambda_a(b)$.  On the
quotient by constants, the derivative seminorm contracts by
$\rho=9/20$.  The block inverse formula therefore gives
\[
 \norm{(I-L_{a,b})^{-1}}
 \le C(1+\abs{1-\lambda_a(b)}^{-1}),
\]
and \cref{lem:phase} proves \eqref{eq:Rb}.  The roof is independent of $a$;
repeated differentiation of the inverse gives words with at most $k+1$
resolvents and the graded Paper-I letters, proving \eqref{eq:Rb-der}.
For $\sigma>0$, the weights acquire $e^{-\sigma\tau_i}$ and the constant block
is uniformly inside the unit disk away from the boundary; at $\sigma=0$ use
the Diophantine estimate.  The quotient remains contracting for small
$\sigma_0$.
\end{proof}

\section{Correct triangular Lorentz geometry}\label{sec:triangle}
Let $\Lambda$ be generated by $(1,0)$ and $(1/2,\sqrt3/2)$, with a disk of
radius $r$ at each lattice point.

\begin{theorem}[Open finite-horizon window]\label{thm:triangle}
If
\[
 {\sqrt3\over4}<r<{1\over2},
\]
the periodic Lorentz gas has disjoint strictly convex scatterers and finite
horizon.  If a $C^2$ normal deformation contains the disk of radius
$r-\delta$, is contained in the disk of radius $r+\delta$, has positive
curvature, and
\[
 r-\delta>{\sqrt3\over4},
 \qquad r+\delta<{1\over2},
\]
then the deformed table also has finite horizon.
\end{theorem}

\begin{proof}
The upper bound is nonoverlap.  An irrational line is dense modulo the lattice
and meets a disk.  A rational direction is parallel to a primitive lattice
vector $v$; the row spacing is
\[
 {\operatorname{area}(\R^2/\Lambda)\over\abs v}
 ={\sqrt3/2\over\abs v}\le{\sqrt3\over2}.
\]
The lower radius bound closes every row corridor.  Compactness gives a uniform
free-flight bound.  The contained disks block corridors after deformation,
and the containing disks remain disjoint.
\end{proof}

\section{Coefficient fields}
For a globally $C^{N,\alpha}$ selector $\Theta(x,p)$ and coupling matrix
$B(x,p)$ define
\[
 A(x,p)={1\over2}B(x,p)\Sigma(\Theta(x,p))B(x,p)^T.
\]

\begin{theorem}[Elliptic coefficient lift]
The drift, covariance, area anomaly, and $A$ are globally $C^{N,\alpha}$.  If
$\Sigma\ge\lambda_\Sigma I$ and $BB^T\ge\lambda_B^2I$, then
\[
 A\ge{1\over2}\lambda_\Sigma\lambda_B^2I.
\]
\end{theorem}

\begin{proof}
Apply the finite-dimensional chain rule and test the quadratic form on
$B^Tz$.
\end{proof}

\section{Conclusion}
The same nonconjugate moving-cut family now has explicit pressure, symmetric
physical diffusion, low-frequency renewal response, and full all-frequency
parameter response.  The corrected triangular theorem supplies the positive
specular finite-horizon benchmark.

\section*{Disclosure and review status}
The author is responsible for all proofs and references.  This controlling
revision awaits a second external specialist review; scripts do not certify
the analysis.

\bibliographystyle{plain}
\bibliography{references}
\end{document}
